\chapter{Sample Doxometric Survey Instruments}
\label{app:survey}

\textbf{Status}: The instruments below are \textbf{pilot items} intended to illustrate the structure of doxometric measurement.
They have not been validated through factor analysis, invariance testing, or external criterion studies.
Before use in research, they require: (a)~exploratory factor analysis to assess dimensionality, (b)~confirmatory factor analysis with holdout samples, (c)~measurement invariance testing across demographic groups, (d)~social desirability bias assessment, and (e)~test-retest reliability estimation.
Where validated instruments exist in the literature (e.g., Kluegel and Smith's stratification beliefs scale, Jost and colleagues' system justification measures), researchers should prefer those.

\section*{Instrument 1: Belief in Selfish Human Nature}

Respondents rate agreement on a 7-point Likert scale (1 = Strongly Disagree, 7 = Strongly Agree).

\begin{enumerate}[nosep]
  \item Most people are primarily motivated by self-interest.
  \item If given the chance, most people would take advantage of others.
  \item Cooperation only happens when people are forced or incentivized.
  \item People who help others usually expect something in return.
  \item Competition is more natural to humans than cooperation.
  \item[\textbf{R6.}] Most people are willing to help a stranger even when there is no reward. \textit{(reverse-coded)}
  \item[\textbf{R7.}] Humans naturally form cooperative communities without being forced. \textit{(reverse-coded)}
  \item[\textbf{R8.}] People generally try to be fair, even when no one is watching. \textit{(reverse-coded)}
\end{enumerate}

\textbf{Scoring}: Reverse-code items R6--R8 before summing. Range 8--56. Higher scores indicate stronger belief in selfish human nature.
Internal consistency (Cronbach's $\alpha$) should be assessed empirically; a threshold of $\alpha > 0.70$ is minimally acceptable.

\textbf{Note}: This instrument conflates beliefs about motivation, norm compliance, and institutional necessity.
Exploratory factor analysis may reveal a multi-factor structure (e.g., ``cynicism about others'' vs.\ ``necessity of enforcement'').
Researchers should test dimensionality before treating it as a single scale.

\section*{Instrument 2: Meritocratic Belief Scale}

\begin{enumerate}[nosep]
  \item People who work hard almost always get ahead.
  \item Success in life is mostly determined by effort and talent.
  \item The wealthy generally deserve their wealth.
  \item Everyone has a roughly equal opportunity to succeed.
  \item[\textbf{R5.}] Success depends more on connections and family background than on individual effort. \textit{(reverse-coded)}
  \item[\textbf{R6.}] The economic system is set up in a way that gives some people unfair advantages. \textit{(reverse-coded)}
\end{enumerate}

\textbf{Scoring}: Reverse-code R5--R6 before summing. Range 6--42. Higher scores indicate stronger meritocratic belief.

\textbf{Note}: This scale folds together beliefs about effort, deservingness, equal opportunity, and poverty attribution.
These may be distinct constructs; pilot testing should assess whether a single-factor model fits the data.

\section*{Instrument 3: Behavioral Vignette (Norm Misperception)}

\begin{quote}
\textit{Imagine you find a wallet containing \$200 and the owner's ID. No one is watching.}
\begin{enumerate}[nosep]
  \item \textit{What would you do?} (Return it / Keep the money / Other)
  \item \textit{What do you think most people would do?} (Return it / Keep the money / Other)
  \item \textit{What percentage of people do you think would return it?} (0--100\%)
\end{enumerate}
\end{quote}

The gap between own behavior (item 1) and perceived norm (items 2--3) measures the \emph{norm misperception gap} (Definition~\ref{def:perceived-norm}).
A large gap (most people say they would return it, but predict others would keep it) indicates doxonomic installation of selfishness as perceived norm.

\section*{General Methodological Notes}

For all instruments:
\begin{itemize}[nosep]
  \item Reverse-code counter-directional items (marked R) before summing.
  \item Report Cronbach's $\alpha$ and factor loadings.
  \item Assess social desirability bias (e.g., include a short social desirability scale and partial out its effects).
  \item For cross-national use, test measurement invariance (configural, metric, scalar) across country groups.
  \item Compare stated-belief scores with behavioral measures to assess divergence (Proposition~\ref{prop:divergence}).
  \item Consider supplementing Likert items with list experiments or endorsement experiments to mitigate response bias on sensitive beliefs.
\end{itemize}
